\section*{Capítulo \ref{ch:g12:subduction-magmatism} --- El magmatismo de las zonas de subducción}

\begin{solution}{exo:g12:subduction-magmatism:1}
(a) un \omterm{def:g12:subduction-magmatism:arc}{arco volcánico}; (b) la \omterm{def:g12:subduction-magmatism:hydrous-melting}{fusión hidratada} (fusión por aporte de agua);
(c) una \omterm{def:g12:subduction-magmatism:granodiorite}{granodiorita}; (d) un \omterm{def:g12:subduction-magmatism:batholith}{batolito}; (e) el \omterm{def:g12:subduction-magmatism:continental-growth}{crecimiento continental}.
\end{solution}

\begin{solution}{exo:g12:subduction-magmatism:2}
(1) Al menos $24/0.25 = \qty{96}{km/h}$, es decir, $\num{96000}/3600 \approx
\qty{27}{m/s}$. (2) La explosión fue al menos $480/96 = 5$ veces más rápida.
\end{solution}

\begin{solution}{exo:g12:subduction-magmatism:3}
(1) $838/1214 \approx 0.690$, alrededor del \qty{69}{\percent}. (2) Del lado
de la placa cabalgante, la que no se hunde. (3) Entre la fosa y el frente, la
\omterm{def:g11:subduction-zones:slab}{losa} aún no es lo bastante profunda: no ha alcanzado la profundidad, de unos
cien kilómetros, a la que libera su agua en un manto lo bastante caliente para
fundirse una vez humedecido.
\end{solution}

\begin{solution}{exo:g12:subduction-magmatism:4}
(1) $45/(45 + 15) = 0.75$: el \qty{75}{\percent} del feldespato es
plagioclasa. (2) Sí: es de grano grueso, contiene más de un
\qty{20}{\percent} de cuarzo, y su fracción de plagioclasa está entre el 65 y
el \qty{90}{\percent}.
\end{solution}

\begin{solution}{exo:g12:subduction-magmatism:5}
Los magmas de arco son más ricos en sílice, así que son más viscosos; y llevan
el agua que los formó, así que son ricos en gas. Una viscosidad alta y un
contenido alto de gas favorecen las dos las erupciones explosivas.
\end{solution}

\begin{solution}{exo:g12:subduction-magmatism:6}
(1) $H = 150 \times \tan\ang{35} \approx 150 \times 0.700 \approx
\qty{105}{km}$. (2) Es la media mundial.
\end{solution}

\begin{solution}{exo:g12:subduction-magmatism:7}
$d = H/\tan\theta$: $105/\tan\ang{20} \approx 105/0.364 \approx
\qty{288}{km}$ a \ang{20}, y $105/\tan\ang{45} = \qty{105}{km}$ a
\ang{45}. Cuanto más suave es la \omterm{def:g11:subduction-zones:slab}{losa}, más lejos de la fosa se alza el arco.
\end{solution}

\begin{solution}{exo:g12:subduction-magmatism:8}
(1) $T = -5.104 \times 3.16^2 + 132.9 \times 3.16 + 1120.7 \approx -51.0 +
420.0 + 1120.7 \approx \qty{1490}{\celsius}$. (2) Unos $1490 - 1000 =
\qty{490}{\celsius}$. (3) En seco, no: \qty{1300}{\celsius} está por debajo de
\qty{1490}{\celsius}. Con agua, sí: está por encima de \qty{1000}{\celsius}.
\end{solution}

\begin{solution}{exo:g12:subduction-magmatism:9}
(1) $2.5/2.7 \approx 0.93$ de año, unos 340 días. (2) La avalancha desplazó
roca antigua que ya existía; los arcos fabrican roca nueva a partir del manto.
Solo se comparan los volúmenes.
\end{solution}

\begin{solution}{exo:g12:subduction-magmatism:10}
(1) $2.7 \times \num{1000000} = \qty{2.7e6}{km^3}$. (2) $\num{9e9}/2.7
\approx \num{3.3e9}$ años, 3300 millones de años. (3) Es menos que la edad de
la Tierra, así que los arcos podrían haber formado toda la \omterm{def:g12:continental-crust:continental-crust}{corteza
continental}; pero la \omterm{def:g11:subduction-zones:subduction}{subducción} también devuelve corteza, así que en realidad
gran parte de la corteza se ha formado y destruido más de una vez.
\end{solution}

\begin{solution}{exo:g12:subduction-magmatism:11}
(1) Los terremotos someros se producen también en la placa cabalgante y en la
falla entre las dos placas; solo los profundos están únicamente en la \omterm{def:g11:subduction-zones:slab}{losa}.
(2) $116 - 80 = \qty{36}{km}$ en \qty{86.4}{km}: $\tan\theta = 36/86.4
\approx 0.417$, $\theta \approx \ang{22.6}$.
\end{solution}

\begin{solution}{exo:g12:subduction-magmatism:12}
En el triángulo rectángulo de la proposición, la \omterm{def:g11:subduction-zones:slab}{losa} es la hipotenusa $L$
y $d$ el cateto contiguo al ángulo $\theta$, así que $\cos\theta = d/L$ y $L =
d/\cos\theta$. Para Japón, $L = 270/\cos\ang{20} \approx 270/0.940 \approx
\qty{287}{km}$. A \qty{83}{mm} al año: $\num{2.87e8}/83 \approx \num{3.5e6}$
años, unos 3.5 millones de años.
\end{solution}

\begin{solution}{exo:g12:subduction-magmatism:13}
(1) $V = \pi r^2 h = \pi \times 0.5^2 \times 8 \approx \qty{6.3}{km^3}$.
(2) Alrededor de $6.3/2.5 \approx 2.5$ veces el volumen de la avalancha.
(3) $P = \rho g z$: a \qty{6}{km}, $2830 \times 9.81 \times 6000 \approx
\qty{1.67e8}{Pa} = \qty{0.17}{GPa}$; a \qty{14}{km}, unos
\qty{3.89e8}{Pa} = \qty{0.39}{GPa}.
\end{solution}

\begin{solution}{exo:g12:subduction-magmatism:14}
(1) $2.7 - 3.2 = \qty{-0.5}{km^3}$ al año, una pérdida de \qty{5e5}{km^3}
por millón de años. (2) $500 \times \num{5e5} = \qty{2.5e8}{km^3}$, es decir,
$\num{2.5e8}/\num{9e9} \approx \qty{2.8}{\percent}$ de la corteza. (3) Cada
ritmo es una estimación, con una incertidumbre mucho mayor que la diferencia
de \qty{0.5}{km^3} al año entre ellos: dentro de los errores, los dos se
equilibran.
\end{solution}

\begin{solution}{exo:g12:subduction-magmatism:15}
(1) $\num{20000} \times 10 = \qty{2e5}{km^3}$. (2) $\num{2e5}/1 =
\num{200000}$ años para todos los arcos continentales del mundo juntos. Un
solo arco recibe solo una pequeña parte de ese magma, así que un \omterm{def:g12:subduction-magmatism:batholith}{batolito}
tarda mucho más: se construye \omterm{def:g11:igneous-rocks:intrusion}{plutón} a \omterm{def:g11:igneous-rocks:intrusion}{plutón}, a lo largo de millones de años.
(3) Al menos los \qty{5}{km} de roca que había sobre su techo.
\end{solution}

\begin{solution}{pb:g12:subduction-magmatism:1}
\textbf{1.} $111.2 \times \cos\ang{39} \approx 111.2 \times 0.777 \approx
\qty{86.4}{km}$.

\textbf{2.} $(144.1 - 141.0) \times 86.4 \approx 3.1 \times 86.4 \approx
\qty{268}{km}$, unos \qty{270}{km}.

\textbf{3.} $(141.0 - 139.7) \times 86.4 \approx \qty{112}{km}$.

\textbf{4.} Sobre la placa cabalgante, la que no se hunde, al oeste de la
fosa: la placa del Pacífico se hunde bajo ella.

\textbf{5.} Allí, la \omterm{def:g11:subduction-zones:slab}{losa} aún no es lo bastante profunda para liberar su agua
en un manto lo bastante caliente para fundirse.

\textbf{6.} $116 - 80 = \qty{36}{km}$.

\textbf{7.} Un grado de longitud, \qty{86.4}{km}.

\textbf{8.} $\tan\theta = 36/86.4 \approx 0.417$, así que $\theta \approx
\ang{22.6}$.

\textbf{9.} La recta sube \qty{36}{km} cada \qty{86.4}{km} hacia el este, así
que alcanza la superficie $98/0.417 \approx \qty{235}{km}$ al este del Iwate,
a $141.0 + 235/86.4 \approx \ang{143.7}$~E: a unos $(144.1 - 143.7) \times 86.4
\approx \qty{35}{km}$ al oeste de la fosa.

\textbf{10.} La \omterm{def:g11:subduction-zones:slab}{losa} no es un plano: entra en la fosa casi horizontal y se
dobla cada vez más hacia abajo a medida que se hunde.

\textbf{11.} $268 \times \tan\ang{20} \approx 268 \times 0.364 \approx
\qty{98}{km}$.

\textbf{12.} Coinciden: \qty{98}{km} por los dos caminos.

\textbf{13.} Ambos están dentro del intervalo mundial, algo menos profundos
que la media de \qty{105}{km}.

\textbf{14.} El Megata está \qty{112}{km} al oeste del Iwate, así que la \omterm{def:g11:subduction-zones:slab}{losa}
bajo él está a unos $98 + 0.417 \times 112 \approx \qty{145}{km}$ de profundidad.

\textbf{15.} En el manto, unos $98 - 35 = \qty{63}{km}$ por debajo de la base
de la corteza.

\textbf{16.} Entró cerca de la dorsal donde se formó la corteza, a partir del
agua de mar; gran parte de ella se guarda en la serpentina.

\textbf{17.} $T = -5.104 \times 3.16^2 + 132.9 \times 3.16 + 1120.7 \approx
\qty{1490}{\celsius}$.

\textbf{18.} En $1490 - 1000 = 490$ a $1490 - 950 = \qty{540}{\celsius}$:
unos \qty{500}{\celsius}.

\textbf{19.} La \omterm{def:g11:subduction-zones:slab}{losa} está varios cientos de grados más fría que el manto que
la rodea, demasiado fría para fundirse incluso con agua: lo que se funde es la
cuña del manto caliente que tiene encima, una vez humedecida.

\textbf{20.} La \omterm{def:g11:subduction-zones:slab}{losa} está a unos \textbf{\qty{100}{km} bajo el frente
volcánico} del noreste de Japón (\qty{98}{km} por los dos métodos).
\end{solution}
